\chapter{Conclusion}

This thesis studied bias-resistant fuzzy TOPSIS group decision-making under coordinated decision-maker contamination. The results show that standard fuzzy TOPSIS and plain bootstrap bagging are vulnerable to targeted manipulation. Reliability-aware methods substantially improve robustness when biased decision makers leave statistically distinguishable rating patterns.

The final recommended proposed methods are M4, M6, and M7. M4 uses reliability-weighted bagging, M6 uses reliability-weighted aggregation, and M7 uses entropy-aware multi-signal reliability with weighted ensemble aggregation. In the main attack-fraction experiments, M4 and M6 preserved the clean target rank through 40\% effective contamination, while M7 preserved the clean target rank across all tested structured contamination levels up to 60\%.

External comparator baselines strengthen the contribution. Median aggregation, trimmed-mean aggregation, MAD-consensus filtering, individual-Borda aggregation, and Huang--Li group-ideal TOPSIS aggregation were unable to match M7 under high structured contamination. Statistical summaries and runtime analysis further support the feasibility and robustness of the approach.

The main limitation is adaptive human-mimic manipulation. The method should be claimed as robust against structured, distinguishable coordinated attacks, not as a universal detector of all bias. Future work may add historical trust scores, supervised reliability learning, interactive consensus feedback, and a public web/API platform for dataset upload and method comparison.
